% discussion_summary.tex
\subsection{Summary of Findings}

This article introduced \texttt{semanticfa}, an R package that runs the
full exploratory factor-analytic workflow---similarity construction,
retention, extraction and diagnostics, interpretation, and refinement---on
scale items' semantic structure, response-free at every stage. The
demonstration ran every exported function on the 50 IPIP Big-Five Factor
Markers and validated results against a conventional factor analysis of
$N = 874{,}434$ respondents, yielding three findings. First, the
instrument's semantic structure is real and largely as intended: a
multi-criterion consensus (parallel analysis 5, Kaiser 6, TEFI 2, EGA 5)
retained five factors that recovered the Big Five domains, and an
independent embedding-depth optimization agreed. Second,
semantic--behavioral coherence is graded by domain: matched Tucker
congruence with the human solution was .92 (Openness), .84
(Conscientiousness), .80 (Neuroticism), .64 (Extraversion), and .52
(Agreeableness), and item-pair similarities correlated $r = .46$ with
raw-response correlations over all 1{,}225 pairs---useful convergence,
informative divergence. Third, encoding item meaning is a substantive
modeling decision: sign-flipped reverse-keyed embeddings became anti-topic
vectors aligning \emph{worst} with keyed human correlations ($r = .137$
vs.\ .434 unflipped), entailment- and cosine-based encodings had
mirror-image strengths (Table~\ref{tab:encodings}), and Monte Carlo
calibration showed conventional fit-index cutoffs do not transfer to
embedding-derived matrices. Refinement proved most practical: eleven
semantically redundant item pairs, a response-free 25-item short form
that \emph{exceeded} the full form in recovering the theoretical partition
(NMI .527 to .810; ARI .402 to .571), and candidate-item vetting that
caught all three planted flaws.
